\section{The canonical divisor}
\label{main:proof-section}

We now apply the preceding results to the canonical divisor. The passage back to the original smooth variety is recorded explicitly, because the numerical-dimension hypothesis concerns a fixed ample twist on \(X\).

\begin{proof}[Proof of Theorem~\ref{main:theorem}]
Apply Proposition~\ref{mmp:approximation} to obtain a finite canonical MMP
\[
X\dashrightarrow Y
\]
and the divisors \(L=K_Y\), \(B_Y\), and \(M_j=L+t_jB_Y\), with generated multiples outside closed subsets \(Z_j\) of codimension at least three. The variety \(Y\) is terminal and therefore smooth in codimension two. These are the standing hypotheses of the numerical results.

Fix a very ample Cartier divisor \(H\) on \(Y\) and a projective resolution \(\pi:\widehat Y\to Y\) that is an isomorphism over the smooth locus. Lemma~\ref{num:growth} gives a fixed Cartier divisor \(P\) on \(\widehat Y\). Choosing one positive integer \(r_0\) for which \(r_0L\) is Cartier, we obtain
\begin{equation}\label{main:quadratic}
h^0(\widehat Y,m\pi^*L+P)
\ \ge\ \frac{m^2}{2}L^2\cdot H^{n-2}+O(m)
\qquad (r_0\mid m).
\end{equation}
It also gives \(L^2\cdot H^{n-2}\ge0\).

Suppose that this intersection number is positive. Choose a smooth projective common resolution \(W\) dominating \(\widehat Y\) and all the models in the finite MMP prefix, with maps
\[
r:W\longrightarrow\widehat Y,\qquad u:W\longrightarrow X.
\]
Proposition~\ref{mmp:approximation} gives
\[
u^*K_X-r^*\pi^*L\ge0.
\]
For every positive multiple \(m\) of \(r_0\), this effective difference gives an injection
\[
H^0(\widehat Y,m\pi^*L+P)
\hookrightarrow H^0(W,mu^*K_X+r^*P).
\]

Set \(P_X=u_*r^*P\). This is a fixed integral Weil divisor, hence Cartier on the smooth \(X\). A rational section \(f\) in the right-hand space satisfies
\[
\operatorname{div}_W(f)+mu^*K_X+r^*P\ge0.
\]
Pushing this inequality forward gives
\[
\operatorname{div}_X(f)+mK_X+P_X\ge0.
\]
Thus there is an injective linear map
\[
H^0(W,mu^*K_X+r^*P)
\hookrightarrow H^0(X,mK_X+P_X).
\]
Choose a single ample Cartier divisor \(A_X\) for which \(A_X-P_X\) has a nonzero section. Multiplication by that fixed section gives
\[
h^0(X,mK_X+A_X)\ge h^0(\widehat Y,m\pi^*L+P).
\]
By \eqref{main:quadratic}, the left-hand side is bounded below by \(cm^2\), for a fixed \(c>0\), along all sufficiently large multiples of \(r_0\). Consequently
\[
\limsup_{m\to\infty}\frac{h^0(X,mK_X+A_X)}{m^2}>0,
\]
contradicting \(\kappa_\sigma(X,K_X)=1\) and the convention \eqref{main:sigma}. It follows that
\[
L^2\cdot H^{n-2}=0.
\]

Proposition~\ref{num:nef} now implies that \(L=K_Y\) is nef, including on curves contained in the singular locus. All other conclusions of Theorem~\ref{main:theorem}, including the effective exceptional comparison \eqref{main:comparison} and its strict support condition, are supplied by the finite MMP in Proposition~\ref{mmp:approximation}.
\end{proof}
